\documentclass[border=6mm]{standalone}
% ==============================================================================
% SETUP.TEX - CONFIGURACIÓN GLOBAL DEL PROYECTO
% ==============================================================================
% Este archivo centraliza todos los paquetes y estilos.
% Debe incluirse en el preámbulo del libro principal y de los archivos standalone.
% \PassOptionsToPackage{gray}{xcolor}
% --- 1. CODIFICACIÓN E IDIOMA ---
% fix-cm habilita las fuentes Computer Modern en tamaños arbitrarios (por debajo
% de 5 pt, entre otros). Debe cargarse antes que fontenc.
\usepackage{fix-cm}
\usepackage[utf8]{inputenc}
\usepackage[T1]{fontenc}
\usepackage[spanish, es-tabla]{babel}  
\usepackage{csquotes} % Recomendado para babel y citas

\usepackage{import}
% mode=image: Usa el PDF pre-compilado, nunca intenta recompilar.
% Debes compilar cada figura manualmente antes de compilar el libro.
\usepackage[mode=image, subpreambles=false]{standalone}

% --- 2. MATEMÁTICAS Y CIENCIAS ---
\usepackage{amsmath, amssymb, amsfonts, mathtools}
\usepackage{enumitem}

% LaTeX trae \DeclareMathSizes{5}{5}{5}{5}, así que a 5 pt (\tiny) los subíndices
% salen del mismo tamaño que la base: en las etiquetas de figuras el M de
% $\omega_M$ queda desproporcionado. Se restablece la proporción habitual.
\DeclareMathSizes{5}{5}{3.7}{3}

% --- 3. DISEÑO Y GEOMETRÍA --- 
% Unificamos los márgenes para todo el documento
\usepackage[left=2.5cm,right=2.5cm,top=3cm,bottom=3cm]{geometry}
\makeatletter
\ifstandalone % Evita que geometry dañe el recorte de las figuras bajándolas 3cm
  \@ifclasswith{standalone}{crop=false}{
    % Es un capítulo/sección: Dejamos que geometry actúe.
  }{
    % Es una figura: Neutralizamos geometry para que no las recorte mal.
    \geometry{pass}
  }
\fi
\makeatother
\usepackage{parskip} % Espaciado entre párrafos sin sangría
\usepackage{float}   % Para usar [H] en figura

% Ajustes de flotantes para evitar espacios verticales exagerados
\renewcommand{\topfraction}{0.95}      % Permite que las figuras ocupen hasta el 95% superior
\renewcommand{\bottomfraction}{0.95}   % Permite que las figuras ocupen hasta el 95% inferior
\renewcommand{\textfraction}{0.05}     % Permite hojas con tan solo un 5% de texto
\renewcommand{\floatpagefraction}{0.8} % Requiere que una página de solo figuras esté 80% llena
\setcounter{topnumber}{4}
\setcounter{bottomnumber}{4}
\setcounter{totalnumber}{10}

% --- 4. GRÁFICOS Y TABLAS ---
\usepackage{graphicx}
\usepackage{booktabs} % Tablas profesionales
\usepackage{caption}  % Control de captions y \captionof
\usepackage{tikz}
\usetikzlibrary{arrows.meta, calc, angles, quotes, babel, backgrounds}
\usepackage{pgfplots}
\usepgfplotslibrary{groupplots}
\usepgfplotslibrary{external}
\usepgfplotslibrary{patchplots}
\pgfplotsset{compat=1.18}
\usepackage[american, straightvoltages]{circuitikz}

% --- 5. COLORES Y CAJAS ---

\usepackage[table]{xcolor}
\usepackage{tcolorbox}

% Definición de colores corporativos/estilo
\definecolor{utnblue}{RGB}{0, 51, 102}
\definecolor{codegreen}{rgb}{0,0.6,0}
\definecolor{codegray}{rgb}{0.5,0.5,0.5}
\definecolor{codepurple}{rgb}{0.58,0,0.82}
\definecolor{backcolour}{rgb}{0.96,0.96,0.96}

% --- 6. ENTORNOS PERSONALIZADOS (CAJAS) ---

% Caja "Resultado" (Azul Corporativo) — Definiciones, fórmulas clave, resultados importantes.
% Uso: \begin{resultado}[Fórmula de Euler] ... \end{resultado}
\newtcolorbox{resultado}[1][]{
    colback=utnblue!5!white,
    colframe=utnblue,
    title=\textbf{#1},
    fonttitle=\bfseries,
    sharp corners,
    boxrule=0.5mm
}

% Caja "Nota" (Gris) — Advertencias, aplicaciones, contexto secundario.
% Uso: \begin{nota}[Cuidado con la calculadora] ... \end{nota}
% Uso: \begin{nota} ... \end{nota}  → título por defecto "Nota"
\newtcolorbox{nota}[1][Nota]{
    colback=codegray!5!white,
    colframe=codegray,
    title=\textbf{#1},
    fonttitle=\bfseries,
    sharp corners,
    boxrule=0.5mm
}

% Caja inline para resaltar un valor y enlazarlo con su otra aparición en una tabla
% o en una cuenta: mismo color, mismo valor.
% Uso: \marcavalor{$4$}  o  \marcavalor[codegreen]{$4$}
\newtcbox{\marcavalor}[1][utnblue]{
    on line,
    boxsep=1pt, left=2pt, right=2pt, top=1pt, bottom=1pt,
    colback=#1!10!white,
    colframe=#1,
    boxrule=0.4pt,
    arc=2pt
}

% --- 7. CÓDIGO FUENTE (LISTINGS) ---
\usepackage{listings}

\lstdefinestyle{miestilo}{
    backgroundcolor=\color{backcolour},   
    commentstyle=\color{codegreen},
    keywordstyle=\color{magenta},
    numberstyle=\tiny\color{codegray},
    stringstyle=\color{codepurple},
    basicstyle=\ttfamily\footnotesize,
    breakatwhitespace=false,         
    breaklines=true,                 
    captionpos=b,                    
    keepspaces=true,                 
    numbers=left,                    
    numbersep=5pt,                  
    showspaces=false,                
    showstringspaces=false,
    showtabs=false,                  
    tabsize=2,
    frame=single,                    
    rulecolor=\color{codegray}
}

\lstset{style=miestilo}
% Soporte para tildes y ñ en bloques de código
\lstset{literate={ñ}{{\~n}}1 {á}{{\'a}}1 {é}{{\'e}}1 {í}{{\'i}}1 {ó}{{\'o}}1 {ú}{{\'u}}1}

% Operadores matemáticos personalizados

% Vector en sentido discreto (lista finita de coordenadas): negrita + flecha.
% Uso: \vect{v}, \vect{e}_k, \vect{0}.
\newcommand{\vect}[1]{\vec{\mathbf{#1}}}

% Fecha de versión y comando para capítulos standalone
\providecommand{\verdate}{\the\year.\ifnum\month<10 0\fi\the\month.\ifnum\day<10 0\fi\the\day}
\newcommand{\chapterversion}{%
    \vspace{-1.5cm}\begin{flushright}\small\textit{\color{codegray}Versión~\verdate}\end{flushright}\vspace{0.5cm}%
}

% --- 8. HIPERVÍNCULOS (DEBE SER EL ÚLTIMO PAQUETE) ---
\usepackage[colorlinks=true, linkcolor=utnblue, urlcolor=utnblue, citecolor=utnblue]{hyperref}

% --- 9. REFERENCIAS CRUZADAS ENTRE CAPÍTULOS STANDALONE ---
% Cuando se compila un capítulo standalone, xr-hyper lee los .aux de los
% otros capítulos (si existen) para resolver \ref{} a labels externos.
% En la compilación del libro completo este bloque no se activa.
\ifstandalone
  \usepackage{xr-hyper}
  \makeatletter
  \newcommand{\xrchap}[1]{%
    \edef\xrc@self{\detokenize{Capitulo#1}}%
    \edef\xrc@job{\jobname}%
    \ifx\xrc@self\xrc@job\else
      \IfFileExists{../#1capitulo/Capitulo#1.aux}%
        {\externaldocument{../#1capitulo/Capitulo#1}}{}%
    \fi
  }
  \makeatother
  \xrchap{01}\xrchap{02}\xrchap{03}\xrchap{04}
  \xrchap{05}\xrchap{06}\xrchap{07}\xrchap{08}
  \xrchap{09}\xrchap{10}\xrchap{11}\xrchap{12}
  \xrchap{13}
\fi
\usetikzlibrary{shadows.blur, shapes.geometric}
% \PassOptionsToPackage{gray}{xcolor}

% Realización directa de una ecuación en diferencias general con retardos unitarios.
% Columnas extremas: cadenas de retardos sobre x[n] (izq.) y sobre y[n] (der.).
% Columna central: cadena de sumadores. El sumador superior entrega y[n].

\begin{document}
% Etiqueta de escalador en caja fija: todos los triángulos quedan del mismo tamaño
% sin importar el ancho del coeficiente.
\newcommand{\glab}[1]{\makebox[19.5pt][c]{$\vphantom{\dfrac{b_0}{a_0}}#1$}}

\begin{tikzpicture}[>=Latex, font=\small,
    gain_r/.style={
        regular polygon, regular polygon sides=3, shape border rotate=30,
        draw=utnblue!60, line width=0.5pt, fill=utnblue!8,
        minimum size=0.8cm, inner sep=0pt, font=\scriptsize,
        blur shadow={shadow blur steps=4, shadow blur extra rounding=0pt,
                     shadow xshift=0.4pt, shadow yshift=-0.6pt, shadow opacity=25}
    },
    gain_l/.style={
        regular polygon, regular polygon sides=3, shape border rotate=210,
        draw=utnblue!60, line width=0.5pt, fill=utnblue!8,
        minimum size=0.8cm, inner sep=0pt, font=\scriptsize,
        blur shadow={shadow blur steps=4, shadow blur extra rounding=0pt,
                     shadow xshift=0.4pt, shadow yshift=-0.6pt, shadow opacity=25}
    },
    delay/.style={
        draw=utnblue!60, line width=0.5pt, fill=utnblue!8,
        minimum width=1.35cm, minimum height=0.8cm, font=\scriptsize,
        blur shadow={shadow blur steps=4, shadow blur extra rounding=0pt,
                     shadow xshift=0.4pt, shadow yshift=-0.6pt, shadow opacity=25}
    },
    sum/.style={
        draw=utnblue!60, line width=0.5pt, circle, minimum size=0.8cm,
        inner sep=0pt, fill=white,
        blur shadow={shadow blur steps=4, shadow blur extra rounding=0pt,
                     shadow xshift=0.3pt, shadow yshift=-0.4pt, shadow opacity=20}
    },
    cross_line/.style={utnblue!40, line width=0.4pt},
    signal/.style={->, thick, utnblue},
    wire/.style={-, thick, utnblue},
    tlab/.style={font=\scriptsize, inner sep=2pt},
]

\def\xL{-5.2}\def\xLg{-2.6}\def\xC{0}\def\xRg{2.6}\def\xR{5.2}
\def\yA{1.4}\def\yB{-0.4}\def\yC{-1.8}\def\yD{-3.2}

% sumador con cruz interna (macro manual)
\newcommand{\suma}[1]{%
  \node[sum] (#1) {};
  \draw[cross_line] ($(#1)+(-0.28,-0.28)$) -- ($(#1)+(0.28,0.28)$);
  \draw[cross_line] ($(#1)+(-0.28,0.28)$) -- ($(#1)+(0.28,-0.28)$);
}

% ---------------- columna central de sumadores ----------------
\begin{scope}
  \path (\xC,\yA) node[sum] (SA){}; \draw[cross_line] (\xC-0.28,\yA-0.28)--(\xC+0.28,\yA+0.28); \draw[cross_line] (\xC-0.28,\yA+0.28)--(\xC+0.28,\yA-0.28);
  \path (\xC,\yB) node[sum] (SB){}; \draw[cross_line] (\xC-0.28,\yB-0.28)--(\xC+0.28,\yB+0.28); \draw[cross_line] (\xC-0.28,\yB+0.28)--(\xC+0.28,\yB-0.28);
  \path (\xC,\yD) node[sum] (SD){}; \draw[cross_line] (\xC-0.28,\yD-0.28)--(\xC+0.28,\yD+0.28); \draw[cross_line] (\xC-0.28,\yD+0.28)--(\xC+0.28,\yD-0.28);
\end{scope}

% espina vertical de sumadores (de abajo hacia arriba)
\draw[signal] (SD.north) -- (\xC,\yC-0.35);
\node at (\xC,\yC) {\large$\vdots$};
\draw[signal] (\xC,\yC+0.35) -- (SB.south);
\draw[signal] (SB.north) -- (SA.south);

% ---------------- columna izquierda: cadena sobre x[n] ----------------
\coordinate (xn)  at (\xL,\yA);
\coordinate (xn1) at (\xL,\yB);
\coordinate (xnM) at (\xL,\yD);
\node[delay] (DLa) at (\xL,{(\yA+\yB)/2}) {Retardo};
\node[font=\footnotesize] at (\xL,\yA+0.55) {$x[n]$};
\draw[signal] (xn) -- (DLa.north);
\draw[signal] (DLa.south) -- (xn1);
\node[font=\footnotesize, anchor=east] at (\xL-0.25,\yB+0.32) {$x[n{-}1]$};
\draw[wire, dashed] (xn1) -- (\xL,\yC+0.2);
\node at (\xL,\yC) {\large$\vdots$};
\draw[wire, dashed] (\xL,\yC-0.2) -- (xnM);
\node[font=\footnotesize] at (\xL,\yD-0.5) {$x[n{-}M]$};

% taps y ganancias b_k/a_0
\node[gain_r] (gb0) at (\xLg,\yA) {\glab{\dfrac{b_0}{a_0}}};
\node[gain_r] (gb1) at (\xLg,\yB) {\glab{\dfrac{b_1}{a_0}}};
\node[gain_r] (gbM) at (\xLg,\yD) {\glab{\dfrac{b_M}{a_0}}};
\fill[utnblue] (xn)  circle (2pt);
\fill[utnblue] (xn1) circle (2pt);
\fill[utnblue] (xnM) circle (2pt);
\draw[signal] (xn)  -- (gb0.west);
\draw[signal] (gb0.east) -- (SA.west);
\draw[signal] (xn1) -- (gb1.west);
\draw[signal] (gb1.east) -- (SB.west);
\draw[signal] (xnM) -- (gbM.west);
\draw[signal] (gbM.east) -- (SD.west);

% ---------------- columna derecha: cadena sobre y[n] ----------------
\coordinate (yn)  at (\xR,\yA);
\coordinate (yn1) at (\xR,\yB);
\coordinate (ynN) at (\xR,\yD);
\node[delay] (DRa) at (\xR,{(\yA+\yB)/2}) {Retardo};
\draw[signal] (yn) -- (DRa.north);
\draw[signal] (DRa.south) -- (yn1);
\node[font=\footnotesize, anchor=west] at (\xR+0.25,\yB+0.32) {$y[n{-}1]$};
\draw[wire, dashed] (yn1) -- (\xR,\yC+0.2);
\node at (\xR,\yC) {\large$\vdots$};
\draw[wire, dashed] (\xR,\yC-0.2) -- (ynN);
\node[font=\footnotesize] at (\xR,\yD-0.5) {$y[n{-}N]$};

% ganancias -a_k/a_0 (apuntan a la izquierda)
\node[gain_l] (ga1) at (\xRg,\yB) {\glab{-\dfrac{a_1}{a_0}}};
\node[gain_l] (gaN) at (\xRg,\yD) {\glab{-\dfrac{a_N}{a_0}}};
\fill[utnblue] (yn1) circle (2pt);
\fill[utnblue] (ynN) circle (2pt);
\draw[signal] (yn1) -- (ga1.east);
\draw[signal] (ga1.west) -- (SB.east);
\draw[signal] (ynN) -- (gaN.east);
\draw[signal] (gaN.west) -- (SD.east);

% ---------------- salida y[n] ----------------
\draw[wire] (SA.east) -- (yn);
\fill[utnblue] (yn) circle (2pt);
\draw[signal] (yn) -- (\xR+1.6,\yA) node[font=\footnotesize, above, pos=0.8, text=black] {$y[n]$};

\end{tikzpicture}
\end{document}
